%% kapa: exam=1 round=29 year=2018 period=2 form=A subject=appraisal_law subject_ko=감정평가관계법규 seq=18 number=18
%% source: Q-Net 감정평가사 기출문제 (cst00309) artlSeq=5202808 file=extracted/1008/5202808/2018년 제29회 감정평가사 1차시험 기출문제/2018년 제29회 감정평가사 1차시험 기출문제/2교시 A형.pdf
%% conversion: kordoc 4.24.51 -> tools/gen1.py (auto); review=proofread-sonnet; flags=none
\begin{question}{18}
부동산 가격공시에 관한 법령상 표준지공시지가의 공시방법에 관한 설명으로 옳은 것은?
\begin{choices}
  \item 국토교통부장관은 매년 표준지공시지가를 중앙부동산가격공시위원회의의 심의를 거쳐 공시하여야 한다. % [원문 그대로]
  \item 표준지공시지가를 공시할 때에는 표준지공시지가의 열람방법을 부동산공시가격시스템에 게시하여야 한다.
  \item 표준지공시지가를 공시할 때에는 표준지공시지가에 대한 이의신청의 기간ㆍ절차 및 방법을 부동산공시가격시스템에 게시하여야 한다.
  \item 국토교통부장관은 표준지공시지가와 표준지공시지가의 열람방법을 표준지소유자에게 개별 통지하여야 한다.
  \item 국토교통부장관은 표준지공시지가를 관보에 공고하고, 그 공고사실을 방송ㆍ신문 등을 통하여 알려야 한다.
\end{choices}
\end{question}
\answer{1}
